% !TEX root = ../main.tex

\section{Seventh Chords (4-Note Chords)}

A 7th chord adds one more third on top of a triad:
\begin{equation*}
    \text{7th Chord} = \text{Root (1)} - \text{3rd (3)} - \text{5th (5)} - \text{7th (7)}
\end{equation*}

\begin{examplebox}{The Octave Cheat Code for 7ths}
Find the root one octave higher, then step downward:
\begin{itemize}
    \item \textbf{Major 7th (Maj7):} Octave Root $- 1$ half-step.
    \item \textbf{Minor 7th (7):} Octave Root $- 2$ half-steps (1 whole step).
\end{itemize}
\end{examplebox}

\section{Common 7th Chord Types}

\begin{itemize}
    \item \textbf{Major 7th (Maj7):} Major Triad $+$ Major 7th ($1 - 3 - 5 - 7$).
    \begin{itemize}
        \item Example: $\text{A}\flat\text{maj7} = A\flat - C - E\flat - G$
    \end{itemize}
    \item \textbf{Dominant 7th (7):} Major Triad $+$ Minor 7th ($1 - 3 - 5 - \flat 7$).
    \begin{itemize}
        \item Naturally occurs on the 5th (Dominant) degree of the major scale.
        \item Has active tension that pulls toward the tonic (home).
        \item Example: $\text{C7} = C - E - G - B\flat$, $\text{G7} = G - B - D - F$
    \end{itemize}
    \item \textbf{Minor 7th (m7):} Minor Triad $+$ Minor 7th ($1 - \flat 3 - 5 - \flat 7$).
    \begin{itemize}
        \item Example: $\text{Gm7} = G - B\flat - D - F$, $\text{Fm7} = F - A\flat - C - E\flat$
    \end{itemize}
\end{itemize}

\section{Half-Diminished Chord \texorpdfstring{($m7\flat 5$ / $\varnothing$)}{(m7b5)}}

\begin{definitionbox}{Half-Diminished Formula}
\textbf{Recipe:} Take a Minor 7th chord ($\text{m7}$) and \textbf{lower the 5th by 1 half-step} ($\flat 5$):
\begin{equation*}
    1 - \flat 3 - \flat 5 - \flat 7
\end{equation*}
\begin{itemize}
    \item \textbf{Why "Half":} The lower triad is diminished ($1 - \flat 3 - \flat 5$), but the 7th is a normal minor 7th ($\flat 7$).
    \item \textbf{Example (from Don't Forget):} F half-diminished $= F - G\sharp - B - D\sharp$
\end{itemize}
\end{definitionbox}
